\documentclass[UTF8,fontset=fandol,11pt,a4paper]{ctexart}
\usepackage[a4paper,margin=1.8cm,headheight=15pt]{geometry}
\usepackage{amsmath,amssymb,mathtools}
\usepackage{booktabs,tabularx,array}
\usepackage{enumitem}
\usepackage[most]{tcolorbox}
\usepackage{tikz}
\usetikzlibrary{arrows.meta,positioning,matrix}
\usepackage{xcolor,fancyhdr,hyperref,microtype,lastpage}

\newcolumntype{Y}{>{\raggedright\arraybackslash}X}
\newcolumntype{L}[1]{>{\raggedright\arraybackslash}p{#1}}
\setlength{\parindent}{2em}
\setlength{\parskip}{0.35em}
\setlength{\emergencystretch}{3em}
\renewcommand{\arraystretch}{1.2}
\setlength{\tabcolsep}{3pt}
\setlist[itemize]{itemsep=0.18em,topsep=0.35em,leftmargin=2em}
\setlist[enumerate]{itemsep=0.18em,topsep=0.35em,leftmargin=2em}

\definecolor{deep}{RGB}{38,58,72}
\definecolor{soft}{RGB}{244,247,248}
\definecolor{greenish}{RGB}{52,91,74}
\definecolor{warnc}{RGB}{136,61,58}
\definecolor{purple}{RGB}{84,72,112}
\definecolor{rulegray}{RGB}{110,116,120}

\hypersetup{colorlinks=true,linkcolor=deep,urlcolor=purple,pdftitle={CO-B01 ISA 语义与数据通路}}
\pagestyle{fancy}
\fancyhf{}
\lhead{\small\color{deep}\textbf{408 · 计组 Bridge CO-B01}}
\rhead{\small\color{deep}ISA Semantic × Datapath}
\cfoot{\small 第 \thepage\ 页 / 共 \pageref{LastPage} 页}
\renewcommand{\headrulewidth}{0.4pt}

\newtcolorbox{corebox}[1]{breakable,colback=soft,colframe=deep,title=\textbf{#1},boxrule=0.8pt,arc=2mm}
\newtcolorbox{methodbox}[1]{breakable,colback=white,colframe=greenish,title=\textbf{#1},boxrule=0.7pt,arc=1.5mm}
\newtcolorbox{warnbox}[1]{breakable,colback=white,colframe=warnc,title=\textbf{#1},boxrule=0.7pt,arc=1.5mm}
\newtcolorbox{examplebox}[1]{breakable,colback=white,colframe=purple,title=\textbf{#1},boxrule=0.7pt,arc=1.5mm}
\newcommand{\kw}[1]{\textbf{\color{deep}#1}}
\newcommand{\code}[1]{\texttt{#1}}

\tikzset{
  co/.style={draw=deep,rounded corners=1.5mm,text width=2.8cm,minimum height=0.95cm,align=center,fill=soft,font=\small,inner sep=3pt},
  state/.style={co,double,double distance=1.2pt,fill=white},
  flow/.style={-{Latex[length=2mm]},thick,draw=deep},
  ref/.style={-{Latex[length=2mm]},dashed,draw=rulegray}
}

\title{\vspace{-1.1cm}\textbf{ISA 语义与数据通路}\\[0.4em]
\large 从软件承诺到硬件实现需求的接口契约}
\author{408 · 计算机组成原理 Bridge CO-B01}
\date{}

\begin{document}
\maketitle

\begin{corebox}{母接口}
CO-02 输出一条指令对软件承诺的状态变化；CO-03 需要的却不是助记符，而是一组可实现、可检查的值与状态需求。两侧通过一个规范化语义包交接：
\[
\boxed{
\text{Instruction Semantic}
\to\text{Architectural Delta Packet}
\to\text{Value Dependency Contract}
\to\text{Datapath / Control Input}}
\]
箭头表示接口翻译，不表示时钟周期。Bridge 只拥有“左侧输出怎样变成右侧合法输入”，不拥有指令定义或具体通路。
\end{corebox}

\begin{methodbox}{本手册的范围}
\begin{itemize}
  \item 左侧提供指令的架构语义：编码、有效地址、访存宽度、下一条 PC 和异常条件。
  \item 右侧需要一组可实现的值依赖、路径、资源调度、控制字、关键路径和提交条件。
  \item 本手册只负责把两者翻译成字段完整的状态差，并做“语义到硬件、硬件到语义”的双向校验。
  \item C/ABI、具体 MUX/总线选择和多指令冒险不在此重复展开；它们分别属于语言调用、单指令通路和流水线问题。
\end{itemize}
\end{methodbox}

\section{术语先行：接口字段不是抽象口号}
\begin{methodbox}{五个接口词}
\begin{itemize}
  \item \kw{语义包：}把一条指令的输入、产生的中间值、允许写入的位置、下一条 PC 和异常条件集中写成的一组字段。
  \item \kw{Read Set：}完成指令语义必须读取的架构状态集合。
  \item \kw{Write Set：}指令允许改变的架构状态集合；不在其中的写使能都是错误副作用。
  \item \kw{Derived Value：}由输入计算出来但尚未提交到架构状态的值，例如有效地址或 ALU 结果。
  \item \kw{接口桥梁：}只翻译两个机制之间的输入/输出格式，不重新定义任一侧的基本规则。
\end{itemize}
\end{methodbox}

\tableofcontents

\section{为什么需要独立接口}

若 CO-02 只给助记符，CO-03 会被迫猜“到底读什么、写什么、异常时哪些效果禁止提交”；若 CO-03 直接从教材图反推语义，又会把一种实现误写成 ISA。接口必须把“软件承诺”转换成与具体微架构无关、却足以生成实现的需求。

这条接口会在三类问题中重复调用：为已有 CPU 增加指令、判断现有通路是否足够、把单指令语义交给流水线/Integration。它因此不是单向 Use 的一句摘要。

\section{左侧输出：架构状态差包}

对指令 $I$ 和初始架构状态 $A$，CO-02 应输出
\[
\Delta_I=
\langle
Read, Derived, Write, NextPC, Memory, Exception
\rangle.
\]

\begin{center}
\small
\begin{tabularx}{\linewidth}{L{2.4cm}Y Y}
\toprule
字段 & 必须回答 & 接口失败的症状\\
\midrule
Read Set & 哪些架构值是语义输入 & 通路少读源或误读目的寄存器\\
Derived Values & EA、算术结果、条件、扩展值等 & 不知道应放哪种功能部件\\
Write Set & 哪些架构位置允许改变 & 多余写使能造成错误副作用\\
Next PC & 顺序、分支、跳转和返回目标 & PC 来源或选择条件不完整\\
Memory Contract & 读/写、宽度、地址、扩展、对齐 & LOAD/STORE 方向或宽度错误\\
Exception Contract & 检测条件、允许提交的效果、重试语义 & faulting 指令留下部分副作用\\
\bottomrule
\end{tabularx}
\end{center}

这个包不包含“第几拍做”“使用几个 ALU”“是否有 MAR/MDR”。这些属于右侧实现选择。

\section{翻译层：从状态差到值依赖契约}

Bridge 依次执行五步：
\begin{enumerate}
  \item \kw{Normalize：}把助记符、字段和寻址方式还原成状态转移，不保留模糊的“执行一下”。
  \item \kw{Expand：}把复合表达式拆为具名中间值，例如 $EA=base+sext(imm)$、$data=M[EA]$。
  \item \kw{Connect：}为每个值写 producer、consumer 和先后依赖。
  \item \kw{Protect：}把 Write Set 与 Exception Contract 转成禁止副作用约束。
  \item \kw{Hand off：}只把值、依赖、允许写入和 memory/exception contract 交给 CO-03。
\end{enumerate}

\begin{center}
\begin{tikzpicture}[node distance=8mm and 7mm]
  \node[state] (isa) {CO-02 Output\\ISA State Delta};
  \node[co,right=of isa] (norm) {Normalize\\六字段语义包};
  \node[co,right=of norm] (dep) {Expand / Connect\\值与依赖};
  \node[state,right=of dep] (cpu) {CO-03 Input\\实现需求};
  \draw[flow] (isa)--(norm);
  \draw[flow] (norm)--(dep);
  \draw[flow] (dep)--(cpu);
  \draw[ref] (cpu.south) to[bend right=25] node[below,font=\scriptsize]{Write Set / Exception 反查} (isa.south);
\end{tikzpicture}
\end{center}

\section{最小例一：ADD}

CO-02 给出
\[
Read=\{R[rs_1],R[rs_2],PC\},\quad
Derived=\{sum,PC_{seq}\},
\]
\[
Write=\{R[rd],PC\},\quad
R'[rd]=R[rs_1]+R[rs_2].
\]
Bridge 将它翻译为两个生产链：寄存器读 $\to$ 加法 $\to$ 目标写回，以及 PC $\to$ 顺序地址 $\to$ PC 写回。CO-03 再决定这些链由一个长周期、多个周期还是其他通路完成。

\begin{warnbox}{接口检查}
若实现打开 MemWrite，Write Set 立即否决；若实现只读取一个源寄存器，Read Set 立即否决。校验不依赖某张标准图。
\end{warnbox}

\section{最小例二：LOAD}

对抽象指令 \code{LOAD rd, imm(rs1)}，左侧语义包为
\[
\begin{aligned}
Read&=\{R[rs1],imm,PC\},\\
EA&=R[rs1]+\operatorname{sext}(imm),\\
data&=M[EA,w],\\
Write&=\{R[rd],PC\}.
\end{aligned}
\]
Bridge 生成依赖
\[
R[rs1],imm\to EA\to memory\ request\to data\to extension\to R[rd].
\]
同时附带：内存必须是读请求、宽度为 $w$、目标写回只能在数据可用且无异常时发生。存储系统何时返回由接口信号表示，不在 Bridge 假定固定一拍。

\begin{examplebox}{为什么这不是 CO-03 的重复正文}
Bridge 只声明 CO-03 必须收到哪些值、依赖和禁止副作用；是否复用 ALU、是否需要临时寄存器、怎样排拍和控制信号为何，是 CO-03 的完整责任。
\end{examplebox}

\section{接口不变量与双向验证}

\begin{itemize}
  \item \kw{Completeness：}每个 ISA 输入都出现在 Read Set 或明确的隐含状态中；每个允许结果都出现在 Write Set。
  \item \kw{No invention：}实现需求不能新增 ISA 未规定的架构副作用。
  \item \kw{Dependency preservation：}consumer 只能使用已经由 producer 产生的值。
  \item \kw{Exception safety：}异常路径不得提交被契约禁止的部分效果。
  \item \kw{Implementation freedom：}接口不限制具体部件数量、拍数或内部寄存器命名。
\end{itemize}

完成设计后进行双向检查：从 ISA 向下确认每个语义效果有实现路径；从通路向上确认每个写使能都能在 Write Set 或 Next PC 中找到授权。

\section{调用时机与停止条件}

\begin{center}
\small
\begin{tabularx}{\linewidth}{L{3.2cm}Y Y}
\toprule
题面信号 & 调用动作 & 何时退出 Bridge\\
\midrule
“增加一条新指令” & 建六字段语义包并展开值依赖 & 进入选部件/加 MUX 时交给 CO-03\\
“判断通路能否支持” & 对照 Read/Derived/Write/NextPC & 开始排拍和算关键路径时退出\\
“写微操作序列” & 先验证序列覆盖语义包 & 具体总线冲突与控制字交给 CO-03\\
“流水线中该指令如何走” & 提供 producer/consumer 与提交条件 & 多指令 need/ready/hazard 交给 CO-04\\
\bottomrule
\end{tabularx}
\end{center}

\section{Anti-Bridge：必须停止的类比}

\begin{itemize}
  \item \kw{ISA 字段 $\neq$ 硬件连线：}字段说明语义来源，不保证直接连到某部件。
  \item \kw{微操作 $\neq$ 指令语义：}微操作是实现步骤，不能反向成为所有机器的 ISA 定义。
  \item \kw{Ready $\neq$ Commit：}值已经产生不表示架构写入已发生。
  \item \kw{C 语句 $\neq$ 唯一指令序列：}编译器和 ISA 可以选择不同表示，只需保持可观察语义。
  \item \kw{Bridge $\neq$ Integration：}本册只拥有可复用接口；完整 LOAD 经 Pipeline/Cache/TLB 的过程属于 CO-I01。
\end{itemize}

\section{一页压缩}

\begin{corebox}{接口四问}
\begin{enumerate}
  \item 指令从哪些架构状态读取？
  \item 必须产生哪些中间值，它们依赖谁？
  \item 最终只允许写哪些架构状态，PC 怎样变化？
  \item 正常/异常路径分别允许哪些效果提交？
\end{enumerate}
回答完四问，CO-02 的语义才能成为 CO-03 可实现、可验证的输入。
\end{corebox}

\end{document}
